\documentclass[10pt,a4paper]{amsart}
\usepackage[margin=19mm]{geometry}
\usepackage{amsmath,amssymb,mathtools}
\usepackage[scr=rsfs,cal=euler]{mathalfa}
\usepackage{xfrac}
\usepackage[unicode,hidelinks,bookmarks=false]{hyperref}
\newcommand{\ie}{i.e.\ }
\newcommand{\cf}{cf.\ }
\newcommand{\resp}{respectively\ }
\newcommand{\Subsection}{\subsection}
\providecommand{\nobreakdash}{\nobreak\hbox{-}}
\setlength{\emergencystretch}{2em}
\multlinegap0pt
\usepackage{booktabs,array}
\usepackage{enumitem}
\newtheorem{theorem}{Theorem}
\newtheorem{proposition}[theorem]{Proposition}
\newtheorem{lemma}[theorem]{Lemma}
\newtheorem{remark}[theorem]{Remark}
\newtheorem{definition}{Definition}[section]
\newcommand{\forb}{\operatorname{forb}}
\newcommand{\Avoid}{\operatorname{Avoid}}
\newcommand{\ex}{\operatorname{ex}}
\newcommand{\cF}{\mathcal F}
\begin{document}
\title[A counterexample to the Anstee-Sali's conjecture]{A counterexample to the Anstee-Sali's conjecture}
\author{Pei Wu}
\date{}
\maketitle
\noindent\emph{Source and licence.} A counterexample to the Anstee-Sali's conjecture. Version arXiv:2608.07646v2 (CC BY 4.0). This material is distributed under the Creative Commons Attribution 4.0 International licence, http://creativecommons.org/licenses/by/4.0/, as recorded in the arXiv OAI licence metadata for this exact version and shown on the abstract page https://arxiv.org/abs/2608.07646v2. Full rights notice: licenses/arxiv-2608.07646-CC-BY-4.0.txt.\par\smallskip
\noindent\textit{Editorial scope.} Counted unit: the original \texttt{proposition} environment \texttt{prop:X4} at source lines 157--162 together with its complete proof at lines 164--206; the delivered document keeps source lines 135--207 so that every referenced label and every citation resolves inside it. Native editable source is the paper's own concatenated TeX; the AMS wrapper is editorial formatting only. No publisher class, upstream script or unrelated artwork is redistributed.
	\section{The six-vertex candidate}
	
	Let
	\[
	V=\{x,y,a,b,c,d\},
	\]
	and define
	\begin{equation}\label{eq:F2}
		\cF_2
		=
		\bigl\{
		\{x,y,a,b\},
		\{x,y,b,c\},
		\{x,y,c,d\},
		\{x,y,d,a\}
		\bigr\}.
	\end{equation}
	Thus $\cF_2$ is the $4$-uniform hypergraph obtained by adjoining the fixed
	core $\{x,y\}$ to every edge of a $4$-cycle on $\{a,b,c,d\}$.
	
	\subsection{\texorpdfstring{Computing $X(\cF_2)$}{Computing X(F2)}}
	

	\begin{proposition}\label{prop:X4}
		For the family $\cF_2$ in \eqref{eq:F2},
		\[
		X(\cF_2)=4.
		\]
	\end{proposition}

	\begin{proof}
		We first prove $X(\cF_2)\le 4$. Consider a four-fold product
		\[
		\mathcal A_1\times\mathcal A_2\times\mathcal A_3\times\mathcal A_4,
		\qquad
		\mathcal A_i\in\{I,I^c,T\},
		\]
		on sufficiently large, pairwise disjoint blocks.
		
		Suppose first that at least one factor is of type $I^c$ or $T$; after
		relabeling, assume it is $\mathcal A_1$. Choose distinct points
		$x,y,a$ in the first block so that the restriction of $\mathcal A_1$ contains
		both $\{x,y\}$ and $\{x,y,a\}$. This is immediate for $I^c$ by omitting $a$
		or a point outside $\{x,y,a\}$, and for $T$ by choosing $x,y,a$ in this order
		along the defining chain. Next choose one point $b,c,d$ from the second,
		third, and fourth blocks, respectively. For each of the standard factors
		$I$, $I^c$, and $T$, the restriction to a single chosen point contains both
		$\varnothing$ and that point. Hence the restriction of the four-fold product
		to $\{x,y,a,b,c,d\}$ contains
		\[
		\{x,y\}\cup B
		\qquad\text{for every }B\subseteq\{a,b,c,d\},
		\]
		and therefore contains $\cF_2$.
		
		It remains to consider the case in which all four factors are of type $I$.
		Choose $x$ from the first block, $y$ from the second block, distinct points
		$a,c$ from the third block, and distinct points $b,d$ from the fourth block.
		The restricted product then contains
		\[
		\{x,y,a,b\},\quad
		\{x,y,a,d\},\quad
		\{x,y,c,b\},\quad
		\{x,y,c,d\},
		\]
		which is exactly a copy of $\cF_2$. Thus every four-fold standard product
		contains $\cF_2$, so $X(\cF_2)\le 4$.
		
		For the reverse inequality, the three-fold product $I\times I\times I$
		avoids $\cF_2$. Every member of this product has size $3$, and restriction
		cannot increase set size, whereas every member of $\cF_2$ has size $4$.
		Therefore $X(\cF_2)>3$, completing the proof.
	\end{proof}
	

\end{document}
